% from ega4-8-fr.tex
\oldpage[IV-3]{44}
De plus, lorsque $g_\lambda$ est une immersion, il en est de même de $g_\mu$
pour $\mu\geq\lambda$.

Il suffit d'appliquer (II, 3.8.4) en remplaçant $\mathcal L$ par
$\mathcal O_X$ et $\mathcal S$ par $\mathcal C[\mathrm T]$, et en tenant
compte de (II, 3.1.7).
\end{lemma}

\begin{proposition}[8.12.3]
\EGAReaderTarget{ega4.c04.prop.8-12-3}{IV.8.12.3-fr}
Soient $Y$ un préschéma quasi-compact et quasi-séparé, $f:X\to Y$ un
morphisme séparé de type fini. Soit $\mathcal B=f_*(\mathcal O_X)$, qui est
une $\mathcal O_Y$-Algèbre quasi-cohérente (I, 9.2.2)~; soit $\mathcal C$ la
fermeture intégrale de $\mathcal O_Y$ dans $\mathcal B$, qui est une
$\mathcal O_Y$-Algèbre quasi-cohérente (II, 6.3.4)~; posons
$Z=\operatorname{Spec}(\mathcal C)$, et soit $g:X\to Z$ le $Y$-morphisme
correspondant à l'injection canonique
$\varphi:\mathcal C\to\mathcal B=f_*(\mathcal O_X)$ (II, 1.2.7). Soit
$(\mathcal C_\lambda)$ la famille filtrante croissante des
sous-$\mathcal O_Y$-Algèbres quasi-cohérentes de type fini de $\mathcal C$
(dont $\mathcal C$ est la réunion ((I, 9.6.6) et (1.7.9))), et, pour tout
$\lambda$, posons $Z_\lambda=\operatorname{Spec}(\mathcal C_\lambda)$, et
soit $g_\lambda:X\to Z_\lambda$ le $Y$-morphisme correspondant à l'injection
$\mathcal C_\lambda\to\mathcal B=f_*(\mathcal O_X)$~; les conditions
suivantes sont équivalentes~:
\begin{enumerate}
\item[\emph{a)}] Il existe une factorisation de $f$ en
\[
X\xrightarrow{f'}Y'\xrightarrow{u}Y,
\]
où $f'$ est une immersion et $u$ un morphisme fini.
\item[\emph{a$'$)}] Il existe une factorisation
$X\xrightarrow{f'}Y'\xrightarrow{u}Y$ de $f$, où $f'$ est une immersion
ouverte et $u$ un morphisme fini.
\item[\emph{b)}] Le morphisme $g:X\to Z$ est une immersion.
\item[\emph{c)}] Il existe $\lambda$ tel que
$g_\lambda:X\to Z_\lambda$ soit une immersion.
\end{enumerate}
De plus, lorsqu'il en est ainsi, $g$ est une immersion ouverte, $g(X)$ est
dense dans $Z$, et il existe $\lambda$ tel que, pour $\mu\geq\lambda$,
$g_\mu$ soit une immersion ouverte.
\end{proposition}

Comme l'homomorphisme $\varphi:\mathcal C\to f_*(\mathcal O_X)$ est
injectif, il résulte de (8.12.1) que si $g$ est une immersion, c'est une
immersion ouverte et $g(X)$ est dense dans $Z$, et il en est de même pour
$g_\lambda$. Le fait que \emph{a)} entraîne \emph{a$'$)} résulte aussi de
(8.12.1), appliqué en remplaçant $Z$ par $Y'$ et $g$ par $f'$ ($Y'$ étant
fini, donc affine au-dessus de $Y$)~: en effet si $Y''$ est l'image fermée de
$X$ par $f'$, $Y''$ est fini au-dessus de $Y$ et $f'$ se factorise en
$X\xrightarrow{f''}Y''\xrightarrow{j}Y'$, où $j$ est l'injection canonique,
et $f''$ est une immersion (I, 4.1.10)~; il résulte donc de (8.12.1) que
$f''$ est une immersion ouverte.

L'équivalence de \emph{b)} et \emph{c)} résulte de (8.12.2) ainsi que le fait
que $g_\lambda$ est alors une immersion pour $\lambda$ assez grand. Il est
clair que \emph{c)} implique \emph{a)}, puisque $Z_\lambda$ est fini
au-dessus de $Y$ (II, 6.3.4 et 6.1.2). Montrons enfin que \emph{a)} implique
\emph{c)}. On a vu ci-dessus qu'on peut supposer que $Y'$ est l'image fermée
de $X$ par $f'$, et alors il résulte de (8.12.1) que si l'on pose
$\mathcal B'=u_*(\mathcal O_{Y'})$, de sorte que $Y'$ s'identifie à
$\operatorname{Spec}(\mathcal B')$, l'homomorphisme
$\varphi':\mathcal B'\to\mathcal B=f_*(\mathcal O_X)$ est \emph{injectif}.
Mais comme par hypothèse $\mathcal B'$ est une $\mathcal O_Y$-Algèbre finie,
elle s'identifie par définition de $\mathcal B$ à une des
sous-$\mathcal O_Y$-Algèbres $\mathcal C_\lambda$, ce qui démontre
\emph{c)}.

Nous dirons qu'un morphisme $f:X\to Y$, où $Y$ est quasi-compact et
quasi-séparé, est \emph{pseudo-fini}, s'il est de type fini et vérifie la
condition \emph{a)} de (8.12.3) (auquel cas il est nécessairement séparé).

\begin{corollary}[8.12.4]
\label{IV.8.12.4-fr}
Soient $Y$ un préschéma quasi-compact et quasi-séparé, $f:X\to Y$ un
morphisme.
